%%% Generated by SAPlugin v10.0 on Sat May 16 12:52:16 CEST 2026

% (optional)
\chapter{Introduction}\label{sec:introduction}
For Software Architecture Part 2 we were asked to add availability and third party data polling to the pms system.


\chapter{QAS Decisions}\label{sec:decisions}



\section{Qual1: Include third party lifestyle tracker dat in clinical model computations (M1)}
\subsection*{Key Decisions}
\begin{adlist}
  \item Including a database in the pms specifically for storing third party patient data

  \item Adding a ThirdPartyApi's component that can be extended to poll third party services for data.

  \item Letting the risk estimation process fetch this third party data and use it for predicting the patient status.
\end{adlist}

\section{Qual2: Communication between patient gateways and PMS backend (Av2)}
\subsection*{Key Decisions}
\begin{adlist}
	\item Introduction of \texttt{SyncManager} and an embedded \texttt{LocalBuffer} module on the Patient Gateway App to abstract network topology and handle a continuous 24-hour network failure gracefully
    \item Implementation of an external \texttt{SMS Gateway} to route emergency notifications over standard GSM signaling when IP-based transport links are non-operational.
	\item Integrating a server-side \texttt{Cronjob} component that acts as a periodic timer loop, auditing liveness states tracking at the \texttt{Patient Gateway Endpoint} to autonomously identify missing transmissions.
    \item Offloading detected communication failures to a dedicated \texttt{AlertManager} component to handle prioritized multi-tiered stakeholder notifications and out-of-band communication channels (such as SMS).
\end{adlist}

% ------------------------------------------------------------
\subsection*{Architectural Decision 1: Adding the cronjob component}

The integration of a decoupled, server-side \texttt{Cronjob} component addresses the fundamental architectural challenge of isolating primary ingestion pathways from orchestration and health-monitoring overhead. In a high-volume telemedicine platform, tracking the live connectivity state of thousands of distributed clients can easily become a performance bottleneck. If the primary network adapters or database access layers are forced to maintain stateful, reactive connection timers for every incoming request, the system risks thread starvation under peak loads, directly threatening overall platform availability.

To mitigate this risk, this decision introduces an asynchronous, pull-based monitoring pattern. The \texttt{Patient Gateway Endpoint} acts as a lightweight, stateless data ingress that simply updates a localized, in-memory metadata table with a \texttt{lastSeen} timestamp whenever a telemetry payload arrives. The \texttt{Cronjob} operates entirely out-of-band as an independent background daemon execution loop, completely decoupled from the main application thread pool. During each execution cycle, it executes a highly optimized, non-blocking batch query against this timestamp index.

By calculating the delta between the current system clock and the expected transmission interval of each active patient profile, the \texttt{Cronjob} can autonomously flag communication dropouts without introducing overhead to the data ingestion pipeline. To preserve high modifiability and strictly adhere to the Single Responsibility Principle, the \texttt{Cronjob} remains entirely agnostic of downstream notification policies or communication mediums. Once a failure state is declared, it hands off a structured exception event to the central \texttt{AlertManager}, which dynamically resolves the stakeholder routing logic.

\subsubsection*{Employed tactics and patterns}
\begin{itemize}
    \item \textbf{Implicit Heartbeat (Availability):} Leveraging existing, periodic scheduled telemetry payloads as connection validation vectors, optimizing client-side battery life and conserving restricted cellular bandwidth.
    \item \textbf{Asynchronous Monitoring Loop (Availability):} Offloading connectivity audits to a distinct background component to insulate core system throughput from monitoring traffic.
\end{itemize}

\subsubsection*{Sensitivity points}
The physical latency of system failure registration is entirely sensitive to the \textit{temporal resolution (polling interval) of the Cronjob scheduler execution loop}. Because the hard upper bound for autonomous backend failure detection is constrained to a 5-minute window, the cron frequency represents a critical system sensitivity point; an overly coarse execution cycle mathematically guarantees a breach of safety constraints, whereas a hyper-aggressive frequency induces severe context-switching overhead on the management node, degrading the performance of adjacent localized services.

\subsubsection*{Relation to other architectural decisions}
This decision establishes a structural dependency on the data contract provided by the \texttt{Patient Gateway Endpoint (AD3)} and coordinates directly with the downstream routing interfaces of the \texttt{AlertManager (AD4)}. When an offline gateway restores its link and flushes its \texttt{LocalBuffer} using an exponential back-off chunking strategy, the endpoint's timestamp index is updated transparently. This allows the subsequent \texttt{Cronjob} execution cycle to self-heal the client's connectivity flag automatically, preventing the generation of redundant failure alerts and protecting telemedicine operators from operational alert fatigue.

\subsubsection*{Trade-off justification}
This architecture enforces a deliberate trade-off between \textbf{System Complexity} and \textbf{Performance Optimization/Modifiability}. Introducing a multi-tiered pipeline (\texttt{Cronjob}, \texttt{Endpoint Ledger}, \texttt{AlertManager}) requires a higher structural overhead and more internal interface specifications than simply embedding connection timeouts inside the active network adapters. However, this trade-off is highly justified within an extra-muros clinical care context; keeping the data ingress entirely stateless ensures that the system can scale seamlessly during mass enrollment phases, while the isolation of the alerting rules guarantees that safety reporting thresholds can be updated over time without requiring redeployments of the primary gateway endpoints.

\subsubsection*{Considered alternatives}
\paragraph{Active Peer-to-Peer Polling from Backend to Client.}
An alternative approach considered was forcing the backend servers to actively broadcast periodic application-layer ping queries to every registered smartphone gateway at 60-second intervals to verify connection health. This alternative was rejected because active background radio polling rapidly drains the edge device's battery, violating the core user experience constraint of low physical invasiveness in the patient's day-to-day life, and scales poorly due to the compounding financial and bandwidth costs of continuous cell-tower polling.
% ------------------------------------------------------------
\subsection*{Architectural Decision 2:
ThirdPartyApi's}

Adding the ThirdPartyApi's component and the googlefit and strava children are present to allow easy implementation of new third party components and a common interface to all.

\subsubsection*{Employed tactics and patterns}
Polymorphism,

\subsubsection*{Sensitivity points}
Here we need to make sure the patient third party token is not leaked they should be securely stored in our database. Implementation of a new third party service should be as fast as possible, a slim parent class would assure this.

\subsubsection*{Relation to other architectural decisions}
A common interface with the third party api's impacted our decissions to use the local database as this allows us to connect these components only once making implementation easier.

\subsubsection*{Trade-off justification}
This makes the application that bit bigger and more complicated but seems well worth it as we are rewarded with better accuracy at a low cost this could also impact latency of the prediction process.

\subsubsection*{Considered alternatives}
\paragraph{Name of alternative 1.}
An alternative could be having a single component to contact all third party apis, however this would lead to a cluttered class once more and more third party apis are added. This parent child structure also gives a clear blueprint for developers to implement.

% ------------------------------------------------------------
\subsection*{Architectural Decision 3:
Patient Gateway Endpoint and the Operator Gateway Endpoint}

This is the entry and exit point for communication between the pms and the patients gateway (their smartphone) or the service operators like doctors, nurses maintenance engineers and such.

\subsubsection*{Employed tactics and patterns}
Single responsiblity,
Role based access control.

\subsubsection*{Sensitivity points}
For these components it is important to make sure they only connect to the components they need. Connecting them to unnecessary components could lead to a user accessing the wrong information.

\subsubsection*{Relation to other architectural decisions}
Making every user type a component is architecturally easy as for every component we have to think what user type needs access to it and connect it with the corresponding endpoint.

\subsubsection*{Trade-off justification}


\subsubsection*{Considered alternatives}
\paragraph{every component just directly contacts the needed user.}
We could have every component that needs access to a user contact him/her directly. This seemed like a bad design to us as connectivity issues are better found if all communication hapens in one component. Our solution allows for less code duplication as there are many common communication functions such as retries and calls to make on failure that can be implemented in the same way for almost all communication calls.

\paragraph{We offer multiple components, each for a certain task to be communicated to a user.}
This seemed like a more viable option compared to the previous alternative but we did not take it as we again have the problem of having to look in multiple places when the communication fails. It also gives a harder time to see what components the gateway can connect to which is something that could be useful for the security team.

% ------------------------------------------------------------
\subsection*{Architectural Decision 4:
Patient DataStorage}

This is our local database for storing the patient data that can't be stored in the hospital information system.

\subsubsection*{Employed tactics and patterns}
Single responsiblity

\subsubsection*{Sensitivity points}
This database stores the incoming data from the third parties that can be contacted for extra patient data. This could include sensitive data so it should be securely implemented and only accessible from the components that need it.

\subsubsection*{Relation to other architectural decisions}
This fits in with the Third party api's as these can directly contact the services and store the data here.

\subsubsection*{Trade-off justification}
The database introduces complexity, however it reduces the time to gather all data necessary for predictions.

\subsubsection*{Considered alternatives}
\paragraph{Storing in the hospital information system (his).}
This would be the best option if we could edit the his to store this data as it moves the security responsibility to the hospital, however it is not possible as we cannot control what can be stored in the his as it depends on the serviced hospital.

% ------------------------------------------------------------
\subsection*{Architectural Decision 5:
HIS Manager}

An abstraction we placed between the hospital information system (his) and the pms for translating calls of the system to the HIS api.

\subsubsection*{Employed tactics and patterns}
Database Abstraction layer

\subsubsection*{Sensitivity points}
This component handles requests to the his database and should do these in as few operations on the his as possible.

\subsubsection*{Relation to other architectural decisions}
This is similar to the endpoint components in that it handles all calls to a external database using the healthAPI.

\subsubsection*{Trade-off justification}
Not including the HIS Manager would require all components to directly connect using the healthAPI, this would be a security issue as our healthAPI key or authentication method would be spread over the whole system. On top of that, it simplifies the transition to another hospital as only this component needs to be rewritten rather than every component using the healtAPI.

\subsubsection*{Considered alternatives}
\paragraph{Polymorphism.}
Just like the third party api's we could implement a parent child structure where we have a common blueprint that is implemented for every hospital. As at the moment we only have one hospital we work with this would only complicate things and not contribute much. However if the pms needs to work with multiple hospital simultaneously this could be a good option.

% (optional)
\chapter{Other decisions}\label{sec:otherdecisions}



\chapter{Discussion}\label{sec:discussion}
